% XeLaTeX; editable table.
\documentclass[UTF8]{ctexart}
\usepackage[a4paper,landscape,margin=18mm]{geometry}
\usepackage{fontspec,booktabs,caption}
\setmainfont{Times New Roman}\pagestyle{empty}
\begin{document}\begin{table}[ht]\centering
{\large Overview of the Evaluated Datasets}\par\medskip
\renewcommand{\arraystretch}{1.25}
\begin{tabular*}{.98\linewidth}{@{\extracolsep{\fill}}lccccccc@{}}
\toprule
 & & & \multicolumn{3}{c}{Tasks} & \multicolumn{2}{c}{Reference Statistics} \\
\cmidrule(lr){4-6}\cmidrule(l){7-8}
Benchmark & Domains & Avg. Tools & User & Injection* & Pairs & Avg. Traj. & Avg. App. \\
\midrule
AgentDojo & 4 & 18.35 & 71 & 35 & 585 & 4.01 & 1.37 \\
ASB & 10 & 3.00 & 51 & 40 & 204 & — & — \\
Synthetic & 6 & 6.00 & 104 & 3 & 104 & — & — \\
\bottomrule\end{tabular*}
\caption{三套数据集在本文实际评测范围内的概览。Domains 为环境或代理类别数；Avg. Tools 为每个用户任务可见的平均工具数，按用户任务加权，ASB 包含每次执行注入的一个攻击工具。User 为不同用户任务数；Injection* 为独立攻击任务或模板数；Pairs 为实际评测的任务—攻击组合数。Avg. Traj. 为每个用户任务的参考工具调用步数；Avg. App. 为参考轨迹涉及的不同应用/服务数。* AgentDojo、ASB 按环境/代理内的攻击任务 ID 去重；Synthetic 按三类攻击模板去重，并实例化为 104 个攻击场景，因此该列不直接比较攻击丰富度。参考统计取 AgentDojo 官方 ground_truth，应用/服务按工具所属客户端模块区分，均按用户任务平均。ASB 与合成数据尚无对应标准参考序列和应用映射，标为“—”。各行均为本文实际评测范围。}
\end{table}\end{document}
